
\documentclass{article}
\usepackage{colortbl}
\usepackage{makecell}
\usepackage{multirow}
\usepackage{supertabular}

\begin{document}

\newcounter{utterance}

\centering \large Interaction Transcript for game `dond', experiment `coop\_it', episode 1 with gpt{-}5{-}mini{-}2025{-}08{-}07{-}no{-}reasoning{-}t1.0.
\vspace{24pt}

{ \footnotesize  \setcounter{utterance}{1}
\setlength{\tabcolsep}{0pt}
\begin{supertabular}{c@{$\;$}|p{.15\linewidth}@{}p{.15\linewidth}p{.15\linewidth}p{.15\linewidth}p{.15\linewidth}p{.15\linewidth}}
    \# & $\;$A & \multicolumn{4}{c}{Game Master} & $\;\:$B\\
    \hline

    \theutterance \stepcounter{utterance}  
    & & \multicolumn{4}{p{0.6\linewidth}}{
        \cellcolor[rgb]{0.9,0.9,0.9}{
            \makecell[{{p{\linewidth}}}]{
                \texttt{\tiny{[P1$\langle$GM]}}
                \texttt{State giocando a un gioco di negoziazione cooperativa in cui dovete accordarvi su come dividere una serie di oggetti tra voi e un altro giocatore.} \\
\\ 
\texttt{Regole:} \\
\texttt{(a) A Lei e all'altro giocatore viene dato un insieme di oggetti. Ognuno di voi riceve anche una funzione di valore segreta, che rappresenta il valore di ciascun tipo di oggetto.} \\
\texttt{(b) Si scambiano messaggi con l'altro giocatore per concordare chi si aggiudica gli oggetti. Potete inviare un massimo di 5 messaggi ciascuno, oppure terminare in anticipo facendo una proposta segreta in qualsiasi momento.} \\
\texttt{(c) A ciascuno di voi viene chiesto di inviare una proposta segreta indicando gli oggetti che desiderate, formattata tra parentesi quadre come segue: "{[}Proposta: <numero> <nome oggetto>, <numero> <nome oggetto>, <...>{]}".} \\
\texttt{(d) Se le vostre proposte sono complementari, cioè ci sono abbastanza oggetti per soddisfare entrambe le proposte, a ciascun giocatore viene assegnato un punteggio basato sulla somma dei valori degli oggetti ricevuti. In caso contrario, entrambi ricevono zero punti.} \\
\texttt{(e) L'obiettivo è massimizzare la somma del proprio punteggio e di quello dell'altro giocatore.} \\
\\ 
\texttt{Cominciamo.} \\
\\ 
\texttt{L'insieme degli oggetti disponibili è:} \\
\\ 
\texttt{2 campane, 1 bambola, 2 fili} \\
\\ 
\texttt{La funzione valore segreta è:} \\
\\ 
\texttt{campana: 1, bambola: 2, filo: 3} \\
\\ 
\texttt{IMPORTANTE: i vostri messaggi, a meno che non si tratti di una proposta segreta, vengono trasmessi direttamente all'altro giocatore, quindi non includete alcuna risposta alle regole o testo di annuncio del vostro messaggio. Per fare una proposta segreta, utilizzate il formato indicato. Non utilizzare delle parentesi quadre quando si comunica all'altro giocatore, altrimenti verrà interpretata come una proposta segreta.} \\
            }
        }
    }
    & \\ \\

    \theutterance \stepcounter{utterance}  
    & \multicolumn{4}{p{0.6\linewidth}}{
        \cellcolor[rgb]{0.9,0.9,0.9}{
            \makecell[{{p{\linewidth}}}]{
                \texttt{\tiny{[P1$\rangle$GM]}}
                \texttt{Ciao! Il mio valore segreto è: campana 1, bambola 2, filo 3. A te cosa valgono gli oggetti? Possiamo puntare a massimizzare la somma: io preferisco soprattutto i fili, poi la bambola, poi le campane.} \\
\\ 
\texttt{Propongo come base: io prendo 2 fili, tu prendi 2 campane e 1 bambola. Se non ti torna, proponi un’alternativa con i tuoi valori.} \\
            }
        }
    }
    & & \\ \\

    \theutterance \stepcounter{utterance}  
    & & & \multicolumn{2}{p{0.3\linewidth}}{
        \cellcolor[rgb]{0.9,0.9,0.9}{
            \makecell[{{p{\linewidth}}}]{
                \texttt{\tiny{[GM$|$GM]}}
                \texttt{continue} \\
            }
        }
    }
    & & \\ \\

    \theutterance \stepcounter{utterance}  
    & & & \multicolumn{2}{p{0.3\linewidth}}{
        \cellcolor[rgb]{0.9,0.9,0.9}{
            \makecell[{{p{\linewidth}}}]{
                \texttt{\tiny{[GM$|$GM]}}
                \texttt{Ciao! Il mio valore segreto è: campana 1, bambola 2, filo 3. A te cosa valgono gli oggetti? Possiamo puntare a massimizzare la somma: io preferisco soprattutto i fili, poi la bambola, poi le campane.} \\
\\ 
\texttt{Propongo come base: io prendo 2 fili, tu prendi 2 campane e 1 bambola. Se non ti torna, proponi un’alternativa con i tuoi valori.} \\
            }
        }
    }
    & & \\ \\

    \theutterance \stepcounter{utterance}  
    & & \multicolumn{4}{p{0.6\linewidth}}{
        \cellcolor[rgb]{0.9,0.9,0.9}{
            \makecell[{{p{\linewidth}}}]{
                \texttt{\tiny{[GM$\rangle$P2]}}
                \texttt{State giocando a un gioco di negoziazione cooperativa in cui dovete accordarvi su come dividere una serie di oggetti tra voi e un altro giocatore.} \\
\\ 
\texttt{Regole:} \\
\texttt{(a) A Lei e all'altro giocatore viene dato un insieme di oggetti. Ognuno di voi riceve anche una funzione di valore segreta, che rappresenta il valore di ciascun tipo di oggetto.} \\
\texttt{(b) Si scambiano messaggi con l'altro giocatore per concordare chi si aggiudica gli oggetti. Potete inviare un massimo di 5 messaggi ciascuno, oppure terminare in anticipo facendo una proposta segreta in qualsiasi momento.} \\
\texttt{(c) A ciascuno di voi viene chiesto di inviare una proposta segreta indicando gli oggetti che desiderate, formattata tra parentesi quadre come segue: "{[}Proposta: <numero> <nome oggetto>, <numero> <nome oggetto>, <...>{]}".} \\
\texttt{(d) Se le vostre proposte sono complementari, cioè ci sono abbastanza oggetti per soddisfare entrambe le proposte, a ciascun giocatore viene assegnato un punteggio basato sulla somma dei valori degli oggetti ricevuti. In caso contrario, entrambi ricevono zero punti.} \\
\texttt{(e) L'obiettivo è massimizzare la somma del proprio punteggio e di quello dell'altro giocatore.} \\
\\ 
\texttt{Cominciamo.} \\
\\ 
\texttt{L'insieme degli oggetti disponibili è:} \\
\\ 
\texttt{2 campane, 1 bambola, 2 fili} \\
\\ 
\texttt{La funzione valore segreta è:} \\
\\ 
\texttt{campana: 0, bambola: 2, filo: 4} \\
\\ 
\texttt{IMPORTANTE: i vostri messaggi, a meno che non si tratti di una proposta segreta, vengono trasmessi direttamente all'altro giocatore, quindi non includete alcuna risposta alle regole o testo di annuncio del vostro messaggio. Per fare una proposta segreta, utilizzate il formato indicato. Non utilizzare delle parentesi quadre quando si comunica all'altro giocatore, altrimenti verrà interpretata come una proposta segreta.} \\
\\ 
\\ 
\texttt{Ciao! Il mio valore segreto è: campana 1, bambola 2, filo 3. A te cosa valgono gli oggetti? Possiamo puntare a massimizzare la somma: io preferisco soprattutto i fili, poi la bambola, poi le campane.} \\
\\ 
\texttt{Propongo come base: io prendo 2 fili, tu prendi 2 campane e 1 bambola. Se non ti torna, proponi un’alternativa con i tuoi valori.} \\
            }
        }
    }
    & \\ \\

    \theutterance \stepcounter{utterance}  
    & & & \multicolumn{4}{p{0.6\linewidth}}{
        \cellcolor[rgb]{0.9,0.9,0.9}{
            \makecell[{{p{\linewidth}}}]{
                \texttt{\tiny{[GM$\langle$P2]}}
                \texttt{Grazie per la trasparenza! Per me i valori sono: campana 0, bambola 2, filo 4. Quindi anch’io valuto molto i fili, poi la bambola, e le campane per me valgono zero.} \\
\\ 
\texttt{La tua proposta massimizza bene la somma: 2 fili a me (valore 8 per me, 6 per te), 2 campane + 1 bambola a te (valore 2+0 per me, 2+2 per te). Totale somma: 10 per me + 7 per te? Aspetta, ricontrollo con i nostri valori:} \\
\texttt{{-} A me 2 fili = 8 (tu li valuti 6)} \\
\texttt{{-} A te 2 campane (tu 2, io 0) + 1 bambola (tu 2, io 2) = 4 per te} \\
\texttt{Somma sociale: 8 + 4 = 12.} \\
\\ 
\texttt{Possiamo forse fare ancora meglio dando anche la bambola a me? Controllo:} \\
\texttt{{-} A me 2 fili + 1 bambola = 8 + 2 = 10} \\
\texttt{{-} A te 2 campane = 2} \\
\texttt{Somma: 12 uguale. Quindi entrambe le divisioni danno somma 12. Preferisco la seconda (io prendo anche la bambola) ma la tua è più equa verso di te.} \\
\\ 
\texttt{Se ti va bene la tua proposta iniziale, accetto volentieri. In alternativa, propongo:} \\
\texttt{Io prendo 2 fili, tu prendi 2 campane e 1 bambola. Se confermi, inviamo entrambi la proposta segreta coerente.} \\
            }
        }
    }
     \\ \\

    \theutterance \stepcounter{utterance}  
    & & & \multicolumn{2}{p{0.3\linewidth}}{
        \cellcolor[rgb]{0.9,0.9,0.9}{
            \makecell[{{p{\linewidth}}}]{
                \texttt{\tiny{[GM$|$GM]}}
                \texttt{continue} \\
            }
        }
    }
    & & \\ \\

    \theutterance \stepcounter{utterance}  
    & & & \multicolumn{2}{p{0.3\linewidth}}{
        \cellcolor[rgb]{0.9,0.9,0.9}{
            \makecell[{{p{\linewidth}}}]{
                \texttt{\tiny{[GM$|$GM]}}
                \texttt{Grazie per la trasparenza! Per me i valori sono: campana 0, bambola 2, filo 4. Quindi anch’io valuto molto i fili, poi la bambola, e le campane per me valgono zero.} \\
\\ 
\texttt{La tua proposta massimizza bene la somma: 2 fili a me (valore 8 per me, 6 per te), 2 campane + 1 bambola a te (valore 2+0 per me, 2+2 per te). Totale somma: 10 per me + 7 per te? Aspetta, ricontrollo con i nostri valori:} \\
\texttt{{-} A me 2 fili = 8 (tu li valuti 6)} \\
\texttt{{-} A te 2 campane (tu 2, io 0) + 1 bambola (tu 2, io 2) = 4 per te} \\
\texttt{Somma sociale: 8 + 4 = 12.} \\
\\ 
\texttt{Possiamo forse fare ancora meglio dando anche la bambola a me? Controllo:} \\
\texttt{{-} A me 2 fili + 1 bambola = 8 + 2 = 10} \\
\texttt{{-} A te 2 campane = 2} \\
\texttt{Somma: 12 uguale. Quindi entrambe le divisioni danno somma 12. Preferisco la seconda (io prendo anche la bambola) ma la tua è più equa verso di te.} \\
\\ 
\texttt{Se ti va bene la tua proposta iniziale, accetto volentieri. In alternativa, propongo:} \\
\texttt{Io prendo 2 fili, tu prendi 2 campane e 1 bambola. Se confermi, inviamo entrambi la proposta segreta coerente.} \\
            }
        }
    }
    & & \\ \\

    \theutterance \stepcounter{utterance}  
    & & \multicolumn{4}{p{0.6\linewidth}}{
        \cellcolor[rgb]{0.9,0.9,0.9}{
            \makecell[{{p{\linewidth}}}]{
                \texttt{\tiny{[P1$\langle$GM]}}
                \texttt{Grazie per la trasparenza! Per me i valori sono: campana 0, bambola 2, filo 4. Quindi anch’io valuto molto i fili, poi la bambola, e le campane per me valgono zero.} \\
\\ 
\texttt{La tua proposta massimizza bene la somma: 2 fili a me (valore 8 per me, 6 per te), 2 campane + 1 bambola a te (valore 2+0 per me, 2+2 per te). Totale somma: 10 per me + 7 per te? Aspetta, ricontrollo con i nostri valori:} \\
\texttt{{-} A me 2 fili = 8 (tu li valuti 6)} \\
\texttt{{-} A te 2 campane (tu 2, io 0) + 1 bambola (tu 2, io 2) = 4 per te} \\
\texttt{Somma sociale: 8 + 4 = 12.} \\
\\ 
\texttt{Possiamo forse fare ancora meglio dando anche la bambola a me? Controllo:} \\
\texttt{{-} A me 2 fili + 1 bambola = 8 + 2 = 10} \\
\texttt{{-} A te 2 campane = 2} \\
\texttt{Somma: 12 uguale. Quindi entrambe le divisioni danno somma 12. Preferisco la seconda (io prendo anche la bambola) ma la tua è più equa verso di te.} \\
\\ 
\texttt{Se ti va bene la tua proposta iniziale, accetto volentieri. In alternativa, propongo:} \\
\texttt{Io prendo 2 fili, tu prendi 2 campane e 1 bambola. Se confermi, inviamo entrambi la proposta segreta coerente.} \\
            }
        }
    }
    & \\ \\

    \theutterance \stepcounter{utterance}  
    & \multicolumn{4}{p{0.6\linewidth}}{
        \cellcolor[rgb]{0.9,0.9,0.9}{
            \makecell[{{p{\linewidth}}}]{
                \texttt{\tiny{[P1$\rangle$GM]}}
                \texttt{Confermo la divisione: tu prendi 2 fili, io prendo 2 campane e 1 bambola. Quando vuoi, inviamo entrambi la proposta segreta coerente.} \\
\\ 
\texttt{{[}Proposta: 2 campane, 1 bambola{]}} \\
            }
        }
    }
    & & \\ \\

    \theutterance \stepcounter{utterance}  
    & & & \multicolumn{2}{p{0.3\linewidth}}{
        \cellcolor[rgb]{0.9,0.9,0.9}{
            \makecell[{{p{\linewidth}}}]{
                \texttt{\tiny{[GM$|$GM]}}
                \texttt{proposal} \\
            }
        }
    }
    & & \\ \\

    \theutterance \stepcounter{utterance}  
    & & & \multicolumn{2}{p{0.3\linewidth}}{
        \cellcolor[rgb]{0.9,0.9,0.9}{
            \makecell[{{p{\linewidth}}}]{
                \texttt{\tiny{[GM$|$GM]}}
                [2, 1, 0]
            }
        }
    }
    & & \\ \\

    \theutterance \stepcounter{utterance}  
    & & \multicolumn{4}{p{0.6\linewidth}}{
        \cellcolor[rgb]{0.9,0.9,0.9}{
            \makecell[{{p{\linewidth}}}]{
                \texttt{\tiny{[GM$\rangle$P2]}}
                \texttt{L'altro giocatore ha presentato una proposta segreta. Ora tocca a Lei presentare una sua proposta.} \\
            }
        }
    }
    & \\ \\

    \theutterance \stepcounter{utterance}  
    & & & \multicolumn{4}{p{0.6\linewidth}}{
        \cellcolor[rgb]{0.9,0.9,0.9}{
            \makecell[{{p{\linewidth}}}]{
                \texttt{\tiny{[GM$\langle$P2]}}
                \texttt{{[}Proposta: 2 fili, 2 campane, 1 bambola{]}} \\
            }
        }
    }
     \\ \\

    \theutterance \stepcounter{utterance}  
    & & & \multicolumn{2}{p{0.3\linewidth}}{
        \cellcolor[rgb]{0.9,0.9,0.9}{
            \makecell[{{p{\linewidth}}}]{
                \texttt{\tiny{[GM$|$GM]}}
                \texttt{proposal} \\
            }
        }
    }
    & & \\ \\

    \theutterance \stepcounter{utterance}  
    & & & \multicolumn{2}{p{0.3\linewidth}}{
        \cellcolor[rgb]{0.9,0.9,0.9}{
            \makecell[{{p{\linewidth}}}]{
                \texttt{\tiny{[GM$|$GM]}}
                [2, 1, 2]
            }
        }
    }
    & & \\ \\

    \theutterance \stepcounter{utterance}  
    & & & \multicolumn{2}{p{0.3\linewidth}}{
        \cellcolor[rgb]{0.9,0.9,0.9}{
            \makecell[{{p{\linewidth}}}]{
                \texttt{\tiny{[GM$|$GM]}}
                [[2, 1, 0], [2, 1, 2]]
            }
        }
    }
    & & \\ \\

\end{supertabular}
}

\end{document}
